\documentclass[11pt,a4paper]{article}
\usepackage[margin=1in]{geometry}
\usepackage{amsmath,amssymb}
\usepackage[utf8]{inputenc}
% \usepackage[T1]{fontenc}  % requires full texlive
\usepackage{hyperref}

\title{A Lean-Verified Crosswalk Between Stochastic Semantics,\\Contracts, and Hoare Logic}
\author{June Kim\\{\small Independent $\cdot$ \texttt{june@june.kim}}}
\date{}

\begin{document}
\maketitle
\thispagestyle{empty}

\noindent
We machine-check in Lean~4 (2,800 lines, zero \texttt{sorry}, standard library only) that behavioral contracts on stochastic morphisms and Hoare triples over $\mathrm{FinStoch}$ are the same structure, proved from two independent angles---one functional, one imperative---whose equivalence is compiled, not just described.

The development models an information-processing pipeline as a Kleisli composition of stochastic kernels. Each stage carries a \emph{contract}: a predicate on outputs preserved under monadic composition. A separate Hoare-logic layer builds Hoare triples and weakest preconditions from scratch, then a bridge theorem shows that \texttt{ContractPreserving\,f\,Q} is equivalent to $\{\top\}\;f\;\{Q\}$---a Hoare triple with trivial precondition. Both proof paths reach the same pipeline correctness result; Lean checks their agreement. A nontriviality counterexample rules out vacuous readings of the specification.

\medskip\noindent
This dual proof assembles a precise crosswalk across communities that have independently formalized overlapping structures under different names:

\begin{itemize}
\item \textbf{Markov categories $\leftrightarrow$ Kleisli($\mathcal{D}$).}
The ambient category is $\mathrm{Kleisli}(\mathcal{D}) = \mathrm{FinStoch}$, the paradigmatic Markov category~\cite{Fritz2020}. The Lean \texttt{Support} typeclass is the possibilistic shadow of Fritz, Perrone, and Rezagholi~\cite{FPR2021}: \texttt{support\_bind} is the monad multiplication for $\mathrm{supp} : V \Rightarrow H$.

\item \textbf{Imperative categories $\leftrightarrow$ contract composition.}
Bonchi, Di~Lavore, Rom\'an, and Staton~\cite{BDRS2025} prove that $\mathrm{Stoch}$ is a posetal imperative category (Corollary~76) and derive Hoare rules as theorems (Theorem~79). The Lean bridge theorem is an independent verification of this identification: behavioral contracts on stochastic morphisms \emph{are} Hoare triples.

\item \textbf{Graded program logics $\leftrightarrow$ handshake conditions.}
The contract predicate is close to a graded Hoare postcondition~\cite{GKOS2021}; the handshake condition between adjacent stages resembles the restricted pointwise order in graded-monadic program logics~\cite{KGSU2026}. These correspondences are plausible rather than exact---the ambient categories differ.

\item \textbf{Categorical cybernetics $\leftrightarrow$ compositional predicates.}
Selection relations on parametrised optics~\cite{CGHR2021} and compositional equilibrium predicates~\cite{GKLF2020} play the same structural role as contracts: behavioral predicates preserved under composition. The ambient structures (optics, games) differ from Kleisli.
\end{itemize}

\noindent
The crosswalk surfaces open problems visible only from the intersection: (1)~the free $\mathbb{N}$-semimodule monad~\cite{KW2021} composed with the support monad morphism would track support cardinality through stochastic composition, connecting Leinster's magnitude~\cite{Leinster2021} to Fritz's possibilistic shadow---this composite has not been constructed; (2)~determinantal point processes have no categorical semantics; (3)~Fritz's categorical entropy has no accompanying achievability or rate-distortion results inside Markov categories.

\medskip\noindent\textbf{Keywords:} Markov categories, Hoare logic, Lean~4, possibilistic shadow, graded monads, imperative categories, weakest precondition, compositional contracts.

\medskip\noindent
The talk is aimed at researchers who work in one of these formalisms but not the others. The Lean development is at \texttt{github.com/kimjune01/natural-framework}. A browsable Rosetta stone with the full crosswalk tables is at \texttt{june.kim/framework-lexicon}.

\begin{thebibliography}{99}
\small

\bibitem{BDRS2025}
F.~Bonchi, E.~Di~Lavore, M.~Rom\'an, S.~Staton.
Program Logics via Distributive Monoidal Categories.
\textit{arXiv:2507.18238}, 2025.

\bibitem{CGHR2021}
M.~Capucci, B.~Gavranovi\'c, J.~Hedges, E.~Rischel.
Towards Foundations of Categorical Cybernetics.
\textit{ACT}, 2021.

\bibitem{Fritz2020}
T.~Fritz.
A synthetic approach to Markov kernels.
\textit{Advances in Mathematics} 370, 2020.

\bibitem{FPR2021}
T.~Fritz, P.~Perrone, S.~Rezagholi.
Three monads on Top and the support as a morphism.
\textit{MSCS} 31(8), 2021.

\bibitem{GKOS2021}
M.~Gaboardi, S.~Katsumata, D.~Orchard, T.~Sato.
Graded Hoare Logic and its Categorical Semantics.
\textit{ESOP}, 2021.

\bibitem{GKLF2020}
N.~Ghani, C.~Kupke, A.~Lambert, F.~Nordvall~Forsberg.
Compositional Game Theory with Mixed Strategies.
\textit{APLAS}, 2020.

\bibitem{KGSU2026}
S.~Kura, M.~Gaboardi, T.~Sekiyama, H.~Unno.
A Category-Theoretic Framework for Dependent Effect Systems.
\textit{ESOP}, 2026.

\bibitem{KW2021}
D.~Kidney, N.~Wu.
Algebras for Weighted Search.
\textit{ICFP}, 2021.

\bibitem{Leinster2021}
T.~Leinster.
\textit{Entropy and Diversity: The Axiomatic Approach}.
Cambridge, 2021.

\bibitem{LCS2025}
J.~Liell-Cock, S.~Staton.
Compositional Imprecise Probability.
\textit{POPL}, 2025.

\bibitem{SK2023}
T.~Sato, S.~Katsumata.
Divergences on Monads for Relational Program Logics.
\textit{MSCS}, 2023.

\end{thebibliography}

\end{document}
